\documentclass[11pt]{article}

\usepackage[a4paper,margin=1in]{geometry}
\usepackage{amsmath,amssymb,bm}
\usepackage{graphicx}
\usepackage[hypertexnames=false]{hyperref}
\usepackage{booktabs}
\usepackage{float}
\usepackage[section]{placeins}

\title{\vspace{-1.5em}
Nonlinear latent manifolds and closure models for projection-based reduced-order models}
\author{S. Ares de Parga, Y. Maday}
\date{\today}

\begin{document}
\maketitle
\vspace{-1.0em}

\begin{abstract}
We compare four ANN-enhanced reduced-order modeling routes under aggressive online compression: (i) closure-based intrusive ROMs (PROM-ANN), (ii) latent-manifold intrusive ROMs (PROM-POD-AE), (iii) non-intrusive direct-coefficient surrogates $(\bm\mu,t)\mapsto\mathbf q_N$, and (iv) non-intrusive latent POD-DL surrogates $(\bm\mu,t)\mapsto\mathbf z\mapsto\mathbf q_N$. For intrusive models, the online state is obtained by time-discrete least-squares Petrov--Galerkin (LSPG) minimization, solved by Gauss--Newton, and accelerated by ECSW hyperreduction. The closure family includes state-conditioned, parameter--time conditioned, and hybrid maps for truncated coefficients. The latent family evolves low-dimensional latent coordinates and reconstructs reduced coefficients through a decoder.

A central methodological choice is ROM-consistent training: learning targets are generated by a linear $n_{\mathrm{tot}}$-dimensional PROM/HPROM solve in the full truncated space, rather than by direct projection of HDM trajectories. We also consider a second enrichment stage built from additional low-cost HPROM trajectories at Latin-hypercube parameter points. This manuscript consolidates the mathematical formulation and experimental protocol for the $250\times250$ 2D parametric inviscid Burgers benchmark and prepares an extension to a shallow-water benchmark.
\end{abstract}

\vspace{-1.0em}

\section{Introduction}
Projection-based model order reduction (PMOR) is a standard framework for reducing high-dimensional parametric models while preserving governing equations \cite{antoulas2005approximation,quarteroni2014reduced}. For nonlinear transport-dominated dynamics, linear subspaces may require large dimensions, which motivates nonlinear manifolds, closure modeling, and latent representations \cite{barnett2023neural,cohen2023nonlinear,aghili2026nonlinear}.

From a machine-learning viewpoint, four families are relevant to this work:
\begin{itemize}
\item intrusive closure-enhanced projection (PROM-ANN) \cite{barnett2023neural},
\item intrusive latent-manifold projection (PROM-POD-AE), in the spirit of hyperreduced nonlinear manifold ROMs \cite{romor2025explicable},
\item non-intrusive direct-coefficient regression maps $(\bm\mu,t)\mapsto\mathbf q_N$,
\item non-intrusive latent POD-DL-ROM maps \cite{fresca2022pod}.
\end{itemize}

\paragraph{Naming convention used in this manuscript.}
When we say \textbf{POD-DL}, we refer to the non-intrusive family in the sense of
Fresca--Manzoni \cite{fresca2022pod}. When we say \textbf{PROM-POD-AE}, we refer
to an intrusive latent-manifold projection route (conceptually aligned with
intrusive nonlinear-manifold HPROM literature such as \cite{romor2025explicable}),
even though those references use different method names.

The goal is a backend-consistent comparison under strong online compression, using the same HDM residual, same time discretization, and same intrusive solver/hyperreduction stack for intrusive variants.

\section{Linear PROM/HPROM Reference and Notation}
\label{sec:linear-rom}

\subsection{POD basis from SVD}
Let $\mathbf u(t;\bm\mu)\in\mathbb R^N$ denote the HDM state with parameters $\bm\mu\in\mathcal D\subset\mathbb R^p$. Define
\begin{equation}
\mathbf u_{\mathrm{ref}}:=\frac{1}{N_s}\sum_{i=1}^{N_s}\mathbf u^i,
\qquad
\mathbf S_c=[\mathbf u^1-\mathbf u_{\mathrm{ref}}\ \cdots\ \mathbf u^{N_s}-\mathbf u_{\mathrm{ref}}].
\end{equation}
Let
\begin{equation}
\mathbf S_c=\mathbf U_{\mathbf S}\mathbf\Sigma_{\mathbf S}\mathbf Y_{\mathbf S}^T
\end{equation}
be its thin SVD, and define the linear ROB
\begin{equation}
\mathbf V_{\mathrm{tot}}\in\mathbb R^{N\times n_{\mathrm{tot}}}
\end{equation}
as the first $n_{\mathrm{tot}}$ left singular vectors of $\mathbf U_{\mathbf S}$.

The linear reduced reconstruction is
\begin{equation}
\tilde{\mathbf u}(t;\bm\mu)=\mathbf u_{\mathrm{ref}}+\mathbf V_{\mathrm{tot}}\mathbf q_N(t;\bm\mu),
\qquad
\mathbf q_N(t;\bm\mu)\in\mathbb R^{n_{\mathrm{tot}}}.
\label{eq:linear-reconstruction}
\end{equation}

\subsection{Linear time-discrete LSPG (PROM)}
Let $\mathbf r^{m+1}(\mathbf u;\bm\mu)\in\mathbb R^N$ denote the HDM time-discrete residual at step $m{+}1$. The linear PROM solve is
\begin{equation}
\mathbf q_N^{m+1}(\bm\mu)
=
\arg\min_{\mathbf x\in\mathbb R^{n_{\mathrm{tot}}}}
\left\|
\mathbf r^{m+1}\!\left(\mathbf u_{\mathrm{ref}}+\mathbf V_{\mathrm{tot}}\mathbf x;\bm\mu\right)
\right\|_2^2.
\label{eq:linear-lspg-min}
\end{equation}

With
\begin{equation}
\mathbf J^{m+1,k}:=\frac{\partial\mathbf r^{m+1}}{\partial\mathbf u},
\qquad
\mathbf T_N:=\mathbf V_{\mathrm{tot}},
\qquad
\mathbf W_N^{m+1,k}:=\mathbf J^{m+1,k}\mathbf T_N,
\end{equation}
the Gauss--Newton increment solves
\begin{equation}
(\mathbf W_N^{m+1,k})^T\mathbf W_N^{m+1,k}\,\Delta\mathbf q_N^{m+1,k}
=
-(\mathbf W_N^{m+1,k})^T
\mathbf r^{m+1}\!\left(\mathbf u_{\mathrm{ref}}+\mathbf V_{\mathrm{tot}}\mathbf q_N^{m+1,k};\bm\mu\right).
\label{eq:linear-gn-update}
\end{equation}

\subsection{Linear HPROM via ECSW}
Assume element-wise residual decomposition
\begin{equation}
\mathbf r^{m+1}(\mathbf u;\bm\mu)
=
\sum_{e_i\in\mathcal E}
\mathbf L_{e_i}^T\,
\mathbf r_{e_i}^{m+1}(\mathbf L_{e_i^+}\mathbf u;\bm\mu).
\end{equation}
ECSW approximation yields
\begin{equation}
\widetilde{\mathbf r}^{m+1}(\mathbf x;\bm\mu)
=
\sum_{e_i\in\widetilde{\mathcal E}}\xi_{e_i}
\mathbf L_{e_i}^T\,
\mathbf r_{e_i}^{m+1}\!\left(
\mathbf L_{e_i^+}\left[\mathbf u_{\mathrm{ref}}+\mathbf V_{\mathrm{tot}}\mathbf x\right];\bm\mu
\right),
\qquad
\xi_{e_i}>0,
\end{equation}
with sampled set $\widetilde{\mathcal E}\subset\mathcal E$ \cite{farhat2015structure,de2023hyper}. The linear HPROM solve is
\begin{equation}
\mathbf q_N^{m+1}(\bm\mu)
=
\arg\min_{\mathbf x\in\mathbb R^{n_{\mathrm{tot}}}}
\left\|
\widetilde{\mathbf r}^{m+1}(\mathbf x;\bm\mu)
\right\|_2^2.
\label{eq:linear-hprom-min}
\end{equation}

\subsection{ROM-consistent training targets}
All learning targets in this work are extracted from trajectories $\mathbf q_N(t;\bm\mu)$ generated by the linear ROM solves \eqref{eq:linear-lspg-min} or \eqref{eq:linear-hprom-min}, depending on the selected backend.

\section{Method I: PROM-ANN Closure Models}
\label{sec:method-ann}

\subsection{Manifold definitions}
For closure models, we introduce the partition
\begin{equation}
\mathbf V_{\mathrm{tot}}=[\mathbf V\ \overline{\mathbf V}],
\qquad
\mathbf q_N=
\begin{bmatrix}
\mathbf q\\
\bar{\mathbf q}
\end{bmatrix},
\qquad
n+\bar n=n_{\mathrm{tot}}.
\end{equation}

PROM-ANN uses retained coordinates $\mathbf q\in\mathbb R^n$ as online unknowns and reconstructs truncated coordinates via learned maps:
\begin{align}
\text{Case 1:}&\quad \bar{\mathbf q}=\mathcal N(\mathbf q),\\
\text{Case 2:}&\quad \bar{\mathbf q}=\mathcal M(\bm\mu,t),\\
\text{Case 3:}&\quad \bar{\mathbf q}=\mathcal H(\mathbf q,\bm\mu,t).
\end{align}

The state approximation is
\begin{equation}
\tilde{\mathbf u}(t;\bm\mu)
=
\mathbf u_{\mathrm{ref}}+
\mathbf V\,\mathbf q(t;\bm\mu)+
\overline{\mathbf V}\,\mathcal F(\mathbf q,\bm\mu,t),
\label{eq:ann-manifold}
\end{equation}
with $\mathcal F\in\{\mathcal N,\mathcal M,\mathcal H\}$.

\subsection{LSPG/Gauss--Newton solve}
At each time step,
\begin{equation}
\mathbf q^{m+1}(\bm\mu)
=
\arg\min_{\mathbf x\in\mathbb R^n}
\left\|
\mathbf r^{m+1}(\tilde{\mathbf u}^{m+1}(\mathbf x;\bm\mu);\bm\mu)
\right\|_2^2.
\label{eq:ann-lspg-min}
\end{equation}

Define
\begin{equation}
\mathbf J^{m+1,k}:=\frac{\partial \mathbf r^{m+1}}{\partial \mathbf u},
\qquad
\mathbf T_q^{m+1,k}:=\frac{\partial \tilde{\mathbf u}^{m+1}}{\partial \mathbf q}
=\mathbf V+\overline{\mathbf V}\frac{\partial\mathcal F}{\partial\mathbf q},
\qquad
\mathbf W_q^{m+1,k}=\mathbf J^{m+1,k}\mathbf T_q^{m+1,k}.
\end{equation}
The Gauss--Newton increment solves
\begin{equation}
(\mathbf W_q^{m+1,k})^T\mathbf W_q^{m+1,k}\,\Delta\mathbf q^{m+1,k}
=
-(\mathbf W_q^{m+1,k})^T
\mathbf r^{m+1}(\tilde{\mathbf u}^{m+1}(\mathbf q^{m+1,k};\bm\mu);\bm\mu).
\end{equation}

\subsection{ECSW hyperreduced form}
PROM-ANN HPROM solves
\begin{equation}
\mathbf q^{m+1}(\bm\mu)
=
\arg\min_{\mathbf x\in\mathbb R^n}
\left\|\widetilde{\mathbf r}^{m+1}(\mathbf x;\bm\mu)\right\|_2^2,
\end{equation}
with test basis $\widetilde{\mathbf W}_q=\widetilde{\mathbf J}\mathbf T_q$.

\subsection{Offline training target}
For Cases 1--3, the supervised targets are extracted from ROM-consistent $\mathbf q_N=[\mathbf q;\bar{\mathbf q}]$ trajectories:
\begin{itemize}
\item Case 1: train $\mathcal N$ from $(\mathbf q,\bar{\mathbf q})$,
\item Case 2: train $\mathcal M$ from $(\bm\mu,t,\bar{\mathbf q})$,
\item Case 3: train $\mathcal H$ from $(\mathbf q,\bm\mu,t,\bar{\mathbf q})$.
\end{itemize}

\section{Method II: PROM-POD-AE Latent Manifold}
\label{sec:method-podae}

\subsection{Autoencoder and latent manifold}
Let
\begin{equation}
\mathbf z=\mathcal E(\mathbf q_N),
\qquad
\widehat{\mathbf q}_N=\mathcal D(\mathbf z),
\qquad
\mathbf z\in\mathbb R^{n_z},\ \ n_z\ll n_{\mathrm{tot}}.
\end{equation}
The intrusive latent manifold is
\begin{equation}
\tilde{\mathbf u}(t;\bm\mu)
=
\mathbf u_{\mathrm{ref}}+\mathbf V_{\mathrm{tot}}\,\mathcal D(\mathbf z(t;\bm\mu)).
\label{eq:podae-manifold}
\end{equation}
This is an intrusive latent-manifold projection approach, in line with intrusive nonlinear-manifold HPROM literature \cite{romor2025explicable}, and is separated from the non-intrusive POD-DL-ROM line \cite{fresca2022pod}.

\subsection{LSPG/Gauss--Newton solve in latent coordinates}
The online unknown is $\mathbf z^{m+1}\in\mathbb R^{n_z}$:
\begin{equation}
\mathbf z^{m+1}(\bm\mu)
=
\arg\min_{\mathbf y\in\mathbb R^{n_z}}
\left\|
\mathbf r^{m+1}(\mathbf u_{\mathrm{ref}}+\mathbf V_{\mathrm{tot}}\mathcal D(\mathbf y);\bm\mu)
\right\|_2^2.
\label{eq:podae-lspg-min}
\end{equation}

Define
\begin{equation}
\mathbf T_z^{m+1,k}:=\frac{\partial\tilde{\mathbf u}^{m+1}}{\partial\mathbf z}
=\mathbf V_{\mathrm{tot}}\frac{\partial\mathcal D}{\partial\mathbf z},
\qquad
\mathbf W_z^{m+1,k}=\mathbf J^{m+1,k}\mathbf T_z^{m+1,k}.
\end{equation}
The Gauss--Newton step solves
\begin{equation}
(\mathbf W_z^{m+1,k})^T\mathbf W_z^{m+1,k}\,\Delta\mathbf z^{m+1,k}
=
-(\mathbf W_z^{m+1,k})^T
\mathbf r^{m+1}(\tilde{\mathbf u}^{m+1}(\mathbf z^{m+1,k};\bm\mu);\bm\mu).
\end{equation}

\subsection{ECSW hyperreduced form}
PROM-POD-AE HPROM solves
\begin{equation}
\mathbf z^{m+1}(\bm\mu)
=
\arg\min_{\mathbf y\in\mathbb R^{n_z}}
\left\|\widetilde{\mathbf r}^{m+1}(\mathbf y;\bm\mu)\right\|_2^2,
\end{equation}
with test basis $\widetilde{\mathbf W}_z=\widetilde{\mathbf J}\mathbf T_z$.

\section{Method III: Data-Driven Direct-Coefficient Baseline}
\label{sec:method-datadriven}

\subsection{Model definition}
The non-intrusive baseline maps exogenous inputs directly to reduced coefficients:
\begin{equation}
\mathbf q_N(t;\bm\mu)\approx\mathcal G_q(\bm\mu,t).
\label{eq:data-driven-map-q}
\end{equation}
This model does not require online residual minimization.

\subsection{Training objective and online cost}
Given ROM-consistent targets $\mathbf q_N^{\mathrm{ref}}(t;\bm\mu)$, one minimizes a supervised loss (e.g., MSE)
\begin{equation}
\mathcal L_{\mathrm{data}}
=\frac{1}{N_{\mathrm{samp}}}\sum_{j=1}^{N_{\mathrm{samp}}}
\left\|\mathcal G_q(\bm\mu_j,t_j)-\mathbf q_{N,j}^{\mathrm{ref}}\right\|_2^2.
\end{equation}
Online prediction is then a direct forward evaluation of $\mathcal G_q$ followed by reconstruction with $\mathbf V_{\mathrm{tot}}$.

\section{Method IV: Data-Driven Latent POD-DL Baseline}
\label{sec:method-datadriven-latent}

\subsection{Model definition}
The second non-intrusive baseline follows the POD-DL family \cite{fresca2022pod}:
\begin{equation}
\mathbf z(t;\bm\mu)\approx\mathcal G_z(\bm\mu,t),\qquad
\mathbf q_N(t;\bm\mu)\approx\mathcal D\!\left(\mathcal G_z(\bm\mu,t)\right),
\label{eq:data-driven-map-z}
\end{equation}
where $\mathcal D$ is a decoder trained offline and $\mathbf z\in\mathbb R^{n_z}$.

\subsection{Training objective and online cost}
This baseline is trained non-intrusively through supervised latent/dynamic losses
(with reconstruction consistency through $\mathcal D$), and evaluated online by a
single forward pass of $\mathcal G_z$ and $\mathcal D$, without residual minimization.

\section{Optional Physics-Regularized Training Note}
\label{sec:prnn-note}

Physics-regularized neural training adds a residual-based term to a supervised objective \cite{chen2021physics,sibuet2025discrete}. In this project, such variants are treated as optional diagnostics (not core baselines) because current results are mixed and significantly more expensive offline.

\section{Burgers Benchmark and Experimental Protocol}
\label{sec:burgers}

\subsection{Governing equations and parameter domain}
We use the 2D parametric inviscid Burgers benchmark from \cite{barnett2023neural},
in conservative vector form:
\begin{equation}
\frac{\partial \mathbf U}{\partial t}
+\frac{\partial \mathbf F(\mathbf U)}{\partial x}
+\frac{\partial \mathbf G(\mathbf U)}{\partial y}
=\mathbf S(x;\bm\mu),
\qquad
\mathbf U=
\begin{bmatrix}
u_x\\ u_y
\end{bmatrix},
\end{equation}
with
\begin{equation}
\mathbf F(\mathbf U)=
\begin{bmatrix}
\tfrac{1}{2}u_x^2\\
\tfrac{1}{2}u_xu_y
\end{bmatrix},
\qquad
\mathbf G(\mathbf U)=
\begin{bmatrix}
\tfrac{1}{2}u_xu_y\\
\tfrac{1}{2}u_y^2
\end{bmatrix},
\qquad
\mathbf S(x;\bm\mu)=
\begin{bmatrix}
0.02\,\exp(\mu_2 x)\\
0
\end{bmatrix}.
\end{equation}
The problem is defined on $\Omega=[0,100]\times[0,100]$ with final time
$T_f=25$ and parameters $\bm\mu=(\mu_1,\mu_2)\in\mathcal D$, where
\begin{equation}
\mathcal D=[4.25,5.50]\times[0.015,0.03].
\end{equation}
This conservative form is the one used by the discrete residual and Jacobian in
the ROM/HPROM pipeline.
The initial/boundary conditions used in our campaign are
\begin{equation}
u_x(0,y,t;\bm\mu)=\mu_1,\qquad
u_x(x,y,0)=1,\qquad
u_y(x,y,0)=1.
\end{equation}

\subsection{Numerical discretization}
The HDM uses a cell-centered finite-volume Godunov discretization on a
$250\times250$ mesh, hence $N=2N_xN_y=125{,}000$. Time integration uses
$500$ uniform steps with $\Delta t=0.05$ (thus $501$ stored states including
$t=0$). The same time-discrete residual is used by linear PROM/HPROM and by all
intrusive ANN-enhanced variants (PROM-ANN and PROM-POD-AE) through
LSPG/Gauss--Newton solves.

\subsection{Training and evaluation protocol}
Stage 1 builds $\mathbf V_{\mathrm{tot}}$ from HDM snapshots at $9$ structured
training parameters (a $3\times3$ tensor grid on $\mathcal D$), giving
$N_s=9\times501=4509$ snapshots. Stage 2 generates ROM-consistent coefficients
$\mathbf q_N(t;\bm\mu)$ by solving the linear full-space reduced model
($n_{\mathrm{tot}}$). Stage 3 trains the learned maps from these reduced
targets.

Evaluation is reported at one verification point and two off-grid points:
\begin{equation}
\bm\mu^{(v)}=(4.875,0.0225),\qquad
\bm\mu^{(1)}=(4.56,0.019),\qquad
\bm\mu^{(2)}=(5.19,0.026).
\end{equation}

An enrichment stage adds $20$ Latin-hypercube points in parameter space.
The enrichment data is generated with reduced solves (no additional HDM runs),
then the same model families are retrained and re-evaluated at the same three
points.

\subsection{Backend-consistency policy}
For this manuscript branch, intrusive comparisons are HPROM-consistent end-to-end:
\begin{itemize}
\item Stage-2 reduced targets generated with linear HPROM in $\mathbf V_{\mathrm{tot}}$,
\item Stage-3 map training from HPROM-consistent datasets,
\item online intrusive evaluation with HPROM backend.
\end{itemize}

\subsection{Reported quantities}
\begin{itemize}
\item trajectory relative errors vs HDM,
\item coefficient-space errors vs linear $n_{\mathrm{tot}}$ reference,
\item per-index absolute/relative coefficient curves,
\item time-resolved coefficient heatmaps,
\item online runtimes and speedups.
\end{itemize}

\section{Shallow-Water Benchmark and Experimental Protocol}
\label{sec:shallowwater}

\subsection{Governing equations and parameterized initial condition}
The second benchmark is a 2D shallow-water model on a flat bottom, in
conservative form:
\begin{equation}
\frac{\partial \mathbf U}{\partial t}
+\frac{\partial \mathbf F(\mathbf U)}{\partial x}
+\frac{\partial \mathbf G(\mathbf U)}{\partial y}=0,
\qquad
\mathbf U=
\begin{bmatrix}
h\\ h u_x\\ h u_y
\end{bmatrix},
\end{equation}
with
\begin{equation}
\mathbf F(\mathbf U)=
\begin{bmatrix}
h u_x\\
h u_x^2+\tfrac{1}{2}gh^2\\
h u_x u_y
\end{bmatrix},
\qquad
\mathbf G(\mathbf U)=
\begin{bmatrix}
h u_y\\
h u_x u_y\\
h u_y^2+\tfrac{1}{2}gh^2
\end{bmatrix}.
\end{equation}

The parametric scenario is a collision of two Gaussian bumps in water depth:
\begin{equation}
h(x,y,0;\bm\mu)=h_0
+\mu_1\exp\!\left(-\frac{\|\mathbf x-\mathbf c_1\|_2^2}{2\sigma^2}\right)
+\mu_2\exp\!\left(-\frac{\|\mathbf x-\mathbf c_2\|_2^2}{2\sigma^2}\right),
\end{equation}
with zero initial velocity components ($u_x=u_y=0$), $h_0=1$, $\sigma=0.045$,
$\mathbf c_1=(0.30,0.50)$, $\mathbf c_2=(0.70,0.50)$, and
\begin{equation}
\bm\mu=(\mu_1,\mu_2)\in[0.06,0.14]\times[0.06,0.14].
\end{equation}
Reflective boundary conditions are imposed on all sides.

\subsection{Numerical setup}
The reference configuration in \texttt{shallowwaters-rom-workbench} uses
$N_x=N_y=600$ on $\Omega=[0,1]\times[0,1]$, so $N=3N_xN_y$ unknowns.
The HDM uses MUSCL reconstruction with HLL/HLLC interface fluxes and an implicit
BDF2 time integrator (BE startup). The default fixed time step is
$\Delta t=7\times10^{-5}$ up to $T_f=0.28$, yielding $4001$ stored states per
parameter value.

\subsection{Sampling protocol aligned with Burgers}
The staged workflow mirrors Burgers:
\begin{itemize}
\item baseline set: $9$ structured parameter points ($3\times3$ grid),
\item enrichment set: $20$ LHS points,
\item Stage 1: POD basis $\mathbf V_{\mathrm{tot}}$ from HDM snapshots,
\item Stage 2: ROM-consistent $\mathbf q_N$ dataset from linear HPROM in
$\mathbf V_{\mathrm{tot}}$,
\item Stage 3: train PROM-ANN (Cases 1--3), PROM-POD-AE, and both data-driven baselines (direct-coefficient and latent POD-DL),
\item online intrusive evaluation with HPROM backend.
\end{itemize}

This keeps the methodology consistent across both benchmarks so that differences
in performance are attributable to the physics and manifold structure, not to
pipeline changes.

\section{Scope for the Current Draft}
\begin{itemize}
\item Core model families: linear HPROM reference, PROM-ANN (Cases 1--3), PROM-POD-AE, data-driven direct-coefficient, and data-driven latent POD-DL-ROM.
\item Baseline/enriched comparison is retained as a controlled data axis.
\item PG/PRNN-style variants are currently kept as optional diagnostic tracks.
\end{itemize}

\bibliographystyle{unsrt}
\bibliography{bibliography}

\end{document}
